\documentclass[11pt]{article}
\usepackage[a4paper,margin=1in]{geometry}
\usepackage{amsmath,amssymb,amsthm}
\usepackage{hyperref}

\newcommand{\dd}{\mathrm{d}}
\newcommand{\E}{\mathbb{E}}
\newcommand{\eps}{\varepsilon}
\newcommand{\cL}{\mathcal{L}}
\newcommand{\cK}{\mathcal{K}}
\newcommand{\cP}{\mathcal{P}}
\renewcommand{\Re}{\operatorname{Re}}
\renewcommand{\Im}{\operatorname{Im}}
\newcommand{\Div}{\operatorname{div}}
\newcommand{\nrm}[2]{\|#1\|_{#2}}

\newtheorem{theorem}{Theorem}
\newtheorem{lemma}[theorem]{Lemma}
\newtheorem{proposition}[theorem]{Proposition}
\newtheorem{corollary}[theorem]{Corollary}
\theoremstyle{remark}
\newtheorem{remark}[theorem]{Remark}

\title{Orbit decorrelation with back-reaction:\\
friction, the Einstein relation, and the kinetic limit\\
in a particle--field model}
\author{Working note}
\date{24 November 2026}

\begin{document}
\maketitle

\section{What is and is not proved}

The note \texttt{orbit\_decorrelation\_theorem.tex} (``the companion note'')
proves orbit decorrelation at $t\sim N^{1/3}$ for a tagged particle in a
\emph{prescribed} Gaussian field. Its list of gaps has, as item~2,
back-reaction: in the particle system the tagged particle polarizes the
medium, which produces a friction of order $1/N$. This note treats that gap
in a model. A particle is coupled to finitely many thermostatted field modes
through a Hamiltonian, so that the particle acts on the field and the Gibbs
measure at temperature $T$ is invariant. A parameter $\kappa\in[0,1]$
multiplies the action of the particle on the field: $\kappa=1$ is the
Hamiltonian model and $\kappa=0$ is the prescribed-field model of the
companion note (for potential fields). Throughout, $\eps=N^{-1/2}$.

\paragraph{Proved.}
\begin{itemize}
  \item Lemma~\ref{lem:gibbs}: the model is well posed and $e^{-H/T}$ is
    annihilated by the adjoint generator. (Invariance of the Gibbs measure
    under the semigroup is then quoted from a standard theorem; it is not
    used anywhere else.)
  \item Theorem~\ref{thm:third} (back-reaction is negligible at $N^{1/3}$):
    for every $\kappa\in[0,1]$, $\E e^{iq\cdot Y_t}$ equals
    $\exp(-\tfrac13\eps^2t^3\,q\cdot D(v_0)q)$ plus an explicit
    error, with the \emph{same} $D$ as in the companion note. At
    $t=sN^{1/3}$ the error is $O(N^{-1/3})$, uniformly in $\kappa$. For
    $\kappa=0$ this improves the rate $N^{-1/6}$ and the range
    $t\ll N^{3/8}$ of the companion note to $N^{-1/3}$ and $t\ll N^{2/5}$
    (for potential fields).
  \item Theorem~\ref{thm:drift} (the friction): the mean velocity satisfies
    $\E v_t-v_0=\eps^2t\,m_\kappa(v_0)+O(\eps^2+\eps^3(t+t^{3/2}))$ with
    \[
      m_\kappa(v)=\Div_vD(v)-\frac{\kappa}{T}\,D(v)\,v ,
    \]
    and $\E Y_t=\tfrac12\eps^2t^2m_\kappa(v_0)+{}$smaller terms, which is
    $O(N^{-1/3})$ at $t\sim N^{1/3}$, as the companion note predicted. The
    first term of $m_\kappa$ is present without back-reaction; the second is
    the polarization drag. Both come from exact algebraic identities
    (Lemma~\ref{lem:id}).
  \item Proposition~\ref{prop:einstein} (Einstein relation): the limiting
    velocity generator is $\nabla\cdot(D\nabla)-\kappa T^{-1}(Dv)\cdot\nabla$.
    For $\kappa=1$ it is symmetric with respect to the Maxwellian
    $e^{-|v|^2/2T}$, which is stationary; the forward equation is
    $\partial_\tau f=\nabla\cdot\bigl(D(\nabla f+vf/T)\bigr)$, the
    ``diffusion plus friction'' structure of the linear Landau and
    Lenard--Balescu operators.
  \item Theorem~\ref{thm:kin} (kinetic time): for $t=\tau\eps^{-2}=\tau N$,
    $|\E\varphi(v_t)-u(\tau,v_0)|\le C\eps$ where $u$ solves the
    Fokker--Planck (backward) equation with the generator above. The theorem
    is stated for a given classical solution $u$ with four bounded
    derivatives; that such a solution exists for $\varphi\in C_b^4$ is a
    classical fact which is quoted, not proved
    (Remark~\ref{rem:reg}).
\end{itemize}

\paragraph{Not proved.}
\begin{itemize}
  \item Nothing is proved about the $N$-particle system. The medium here is
    a finite set of field modes, each damped at a rate $\gamma_k>0$ and
    driven by white noise. The damping stands in for Landau damping and is
    put in by hand; the resonances are Lorentzians of width $\gamma_k$, not
    $\delta(k\cdot(v-v_*))$ divided by a dielectric function. The constants
    grow as $\gamma_k\to0$.
  \item The coefficient of the $N^{-1/3}$ correction to the characteristic
    function of $Y$ at $t\sim N^{1/3}$ is not identified. The mean
    displacement is (Theorem~\ref{thm:drift}); the characteristic function
    has other terms of the same order (Remark~\ref{rem:next}).
  \item The field is still Markovian and Gaussian apart from the source;
    gaps 1 and 3--6 of the companion note are untouched.
\end{itemize}

\paragraph{Prior art.}
None of the methods is new. A particle coupled linearly to a bath of
oscillators, with the bath eliminated to give friction and noise related by
fluctuation--dissipation, is the Ford--Kac--Mazur, Caldeira--Leggett and
Kac--Zwanzig construction \cite{FKM,CL,Zwanzig}; the Einstein relation in that
setting is Kubo's \cite{Kubo}. The model here differs in that the coupling
$e^{ik\cdot x}$ is nonlinear in the particle position and the bath modes
carry their own thermostat, so the reduced equation is not a linear
generalized Langevin equation. The proofs use perturbed test functions with
explicit correctors for a fast Markov process
\cite{PSV,PavliotisStuart}; the limit theorem for a particle in a weak
random force field is Kesten--Papanicolaou \cite{KP}. In plasma test-particle
theory the two parts of the drift are standard: the drag from the
polarization wake of the test particle and the drift from the velocity
dependence of the diffusion tensor, with the Lenard--Balescu operator as
their sum \cite{RR,TH,Hubbard,Lenard,Balescu}. What is done here is to prove
the corresponding statements, with rates, in a model that keeps the
resonance structure $\omega_k=k\cdot v$ and is simple enough for every
corrector to be written down. The references are cited from memory and were
not checked against the sources.

\section{The model}

\subsection{Hamiltonian and thermostat}

Fix a finite set $K\subset\mathbb{R}^d\setminus\{0\}$ and for $k\in K$
numbers $c_k\in\mathbb{C}$, $\gamma_k>0$, $\omega_k\in\mathbb{R}$; fix $T>0$
and $\eps\in(0,1]$. The state is $(x,v,a)$ with $x,v\in\mathbb{R}^d$ and
$a=(a_k)_{k\in K}\in\mathbb{C}^K$. Let
\begin{equation}
  H(x,v,a)=\tfrac12|v|^2+T\sum_{k\in K}|a_k|^2
    +\eps\sum_{k\in K}\Re\bigl(c_k\,a_k\,e^{ik\cdot x}\bigr).
  \label{eq:H}
\end{equation}
Write $a_k=a_k'+ia_k''$ and
$\partial_{\bar a_k}=\tfrac12(\partial_{a_k'}+i\partial_{a_k''})$. Let
$(B_k)$ be independent complex Brownian motions with $\E|B_k(t)|^2=t$,
$\E B_k(t)^2=0$. The dynamics is
\begin{equation}
\begin{aligned}
  \dd x&=v\,\dd t,\qquad
  \dd v=-\nabla_xH\,\dd t
     =\eps\sum_k\Re\bigl(f_k\,a_k\,e^{ik\cdot x}\bigr)\dd t,
     \qquad f_k=-i\,c_k\,k,\\
  \dd a_k&=-\frac{\gamma_k+i\omega_k}{T}\,\partial_{\bar a_k}H\,\dd t
     +\sqrt{2\gamma_k}\,\dd B_k
   =-(\gamma_k+i\omega_k)\Bigl(a_k+\frac{\eps\,\bar c_k}{2T}\,e^{-ik\cdot x}\Bigr)\dd t
     +\sqrt{2\gamma_k}\,\dd B_k .
\end{aligned}
\label{eq:model}
\end{equation}
The part with $i\omega_k$ is Hamiltonian; the part with $\gamma_k$, together
with the noise, is a Langevin thermostat acting on the field only. With the
source term $\eps\bar c_ke^{-ik\cdot x}/2T$ dropped, $a_k$ is the stationary
Ornstein--Uhlenbeck mode of the companion note and
$F=\sum_k\Re(f_ka_ke^{ik\cdot x})$ its potential field.

\begin{lemma}\label{lem:gibbs}
\emph{(i)} For every initial condition, \eqref{eq:model} has a unique global
strong solution. \emph{(ii)} Let $\cL_H$ be the generator of
\eqref{eq:model} and $\rho=e^{-H/T}$. Then $\cL_H^*\rho=0$ pointwise, where
$\cL_H^*$ is the formal adjoint with respect to Lebesgue measure on
$\mathbb{R}^d\times\mathbb{R}^d\times\mathbb{C}^K$. \emph{(iii)} With
$\tilde a_k=a_k+\eps\bar c_ke^{-ik\cdot x}/2T$,
\begin{equation}
  H=\tfrac12|v|^2+T\sum_k|\tilde a_k|^2-\frac{\eps^2}{4T}\sum_k|c_k|^2 ,
  \label{eq:Hdressed}
\end{equation}
so under $\rho$ the variables $x$, $v$, $(\tilde a_k)$ are independent: $x$
has Lebesgue measure, $v$ is $N(0,T\,\mathrm{Id})$, and the $\tilde a_k$ are
standard complex Gaussians ($\E|\tilde a_k|^2=1$, $\E\tilde a_k^2=0$).
\end{lemma}

\begin{proof}
(i) Let $\hat a_k$ solve
$\dd\hat a_k=-(\gamma_k+i\omega_k)\hat a_k\dd t+\sqrt{2\gamma_k}\dd B_k$,
$\hat a_k(0)=a_k(0)$. For any continuous path $x$, the third line of
\eqref{eq:model} is linear in $a_k$ and gives
\[
  a_k(t)=\hat a_k(t)-\frac{\eps\bar c_k}{2T}(\gamma_k+i\omega_k)
    \int_0^te^{-(\gamma_k+i\omega_k)(t-s)}e^{-ik\cdot x_s}\,\dd s ,
  \qquad
  |a_k(t)-\hat a_k(t)|\le\frac{\eps|c_k|\,|\gamma_k+i\omega_k|}{2T\gamma_k}.
\]
The coefficients of \eqref{eq:model} are smooth and the noise is additive,
so there is a unique strong solution up to a possible explosion time. On
any interval of existence the display bounds $|a_k|$ by
$|\hat a_k|+C$, hence the force on the particle by
$\eps\sum_k|f_k|(|\hat a_k(t)|+C)$, which is locally bounded in $t$; so
$v$ and $x$ stay bounded on bounded time intervals and there is no
explosion.

(ii) In the real coordinates $(a_k',a_k'')$ the drift of $a_k$ is
\[
  \frac{\omega_k}{2T}\bigl(\partial_{a_k''}H,\,-\partial_{a_k'}H\bigr)
  -\frac{\gamma_k}{2T}\bigl(\partial_{a_k'}H,\,\partial_{a_k''}H\bigr),
\]
and $\sqrt{2\gamma_k}B_k$ has independent real and imaginary parts of
variance $\gamma_kt$. So $\cL_H=X+\sum_k\cL_k$ with
\[
  X=v\cdot\nabla_x-\nabla_xH\cdot\nabla_v
    +\sum_k\frac{\omega_k}{2T}\bigl(\partial_{a_k''}H\,\partial_{a_k'}
      -\partial_{a_k'}H\,\partial_{a_k''}\bigr),
  \qquad
  \cL_k=\frac{\gamma_k}{2}\,e^{H/T}\,\nabla_{a_k}\!\cdot\!
     \bigl(e^{-H/T}\nabla_{a_k}\,\cdot\,\bigr),
\]
$\nabla_{a_k}=(\partial_{a_k'},\partial_{a_k''})$. The vector field $X$ is
divergence free and $XH=0$, so $X^*\rho=-X\rho=0$. For each $k$,
$\cL_k^*\rho=\frac{\gamma_k}{2}\nabla_{a_k}\cdot
(\rho\,\nabla_{a_k}(\rho^{-1}\rho))=0$. (The identity $\cL_H^*\rho=0$ was
also checked symbolically for $d=1$ and one mode.)

(iii) $T|\tilde a_k|^2=T|a_k|^2+\eps\Re(c_ka_ke^{ik\cdot x})
+\eps^2|c_k|^2/4T$.
\end{proof}

\begin{remark}[Invariance]\label{rem:inv}
If $K\subset2\pi\mathbb{Z}^d$ the dynamics descends to
$x\in\mathbb{T}^d$ and $\rho$ normalizes to a probability measure, the Gibbs
measure. That it is invariant under the semigroup, and not only
infinitesimally, follows from (i), (ii) and a standard argument: the
martingale problem for $\cL_H$ on $C_c^\infty$ is well posed because the
coefficients are locally Lipschitz and there is no explosion
\cite{SV}, and then $\int\cL_Hf\,\rho=0$ for $f\in C_c^\infty$ implies
invariance by Echeverr\'ia's theorem \cite[Ch.~4, Thm~9.17]{EK}. This is
quoted, not proved here (the theorem number is from memory), and nothing
below depends on it. The theorems below start the particle at a fixed
$(x_0,v_0)$, not from the Gibbs measure.
\end{remark}

\subsection{Dressed variables and the interpolating family}

By \eqref{eq:Hdressed} the force on the particle is
$-\nabla_xH=\eps\sum_k\Re(f_k\tilde a_ke^{ik\cdot x})$: the self-force
$\eps^2\Re(ik|c_k|^2)/2T$ vanishes. Since
$\dd(e^{-ik\cdot x})=-i(k\cdot v)e^{-ik\cdot x}\dd t$,
\[
  \dd\tilde a_k=\Bigl[-(\gamma_k+i\omega_k)\tilde a_k
     -i\eps\,(k\cdot v)\,\frac{\bar c_k}{2T}\,e^{-ik\cdot x}\Bigr]\dd t
     +\sqrt{2\gamma_k}\,\dd B_k .
\]
In the dressed variable the source is proportional to $k\cdot v$: a particle
at rest does not excite the field. For $\kappa\in[0,1]$ consider
\begin{equation}
\begin{aligned}
  \dd x&=v\,\dd t,\qquad
  \dd v=\eps\sum_k\Re\bigl(f_k\,\tilde a_k\,e^{ik\cdot x}\bigr)\dd t,\\
  \dd\tilde a_k&=\Bigl[-(\gamma_k+i\omega_k)\tilde a_k
     -i\eps\kappa\,(k\cdot v)\,\frac{\bar c_k}{2T}\,e^{-ik\cdot x}\Bigr]\dd t
     +\sqrt{2\gamma_k}\,\dd B_k .
\end{aligned}
\label{eq:kappa}
\end{equation}
For $\kappa=1$ this is \eqref{eq:model}. For $\kappa=0$ it is the
prescribed-field model of the companion note with $f_k=-ic_kk$. Only
$\kappa\in\{0,1\}$ has a physical meaning; intermediate values are used to
keep track of where the back-reaction enters.

\subsection{Reduction and initial data}

Let $z_k=\tilde a_k\,e^{ik\cdot x}$, the amplitude of mode $k$ seen by the
particle, and
\begin{equation}
  \lambda_k(v)=\gamma_k+i(\omega_k-k\cdot v),\qquad
  s_k(v)=-i\,(k\cdot v)\,\frac{\bar c_k}{2T},\qquad
  G(z)=\sum_k\Re(f_kz_k).
\end{equation}
By It\^o's formula ($x$ has finite variation),
\begin{equation}
  \dd v=\eps\,G(z)\,\dd t,\qquad
  \dd z_k=\bigl[-\lambda_k(v)\,z_k+\eps\kappa\,s_k(v)\bigr]\dd t
     +\sqrt{2\gamma_k}\,\dd W_k,
  \label{eq:vz}
\end{equation}
where $W_k=\int e^{ik\cdot x}\dd B_k$ are again independent complex Brownian
motions with $\E|W_k(t)|^2=t$, $\E W_k(t)^2=0$ (L\'evy's characterization:
$\dd[W_k,\bar W_l]=\delta_{kl}\dd t$, $\dd[W_k,W_l]=0$). So $(v,z)$ is
itself a diffusion: the position has been eliminated. Let
$Y_t=x_t-x_0-v_0t$ and $w_t=v_t-v_0$, so $\dot Y=w$. The generator of
$(Y,v,z)$ is
\begin{equation}
  \cL=(v-v_0)\cdot\nabla_Y+\eps\,G\cdot\nabla_v+A^v+\eps\kappa\,S,
  \label{eq:gen}
\end{equation}
\[
  A^v=\sum_k\Bigl(-\lambda_k(v)z_k\partial_{z_k}
     -\bar\lambda_k(v)\bar z_k\partial_{\bar z_k}
     +2\gamma_k\partial_{z_k}\partial_{\bar z_k}\Bigr),
  \qquad
  S=\sum_k\bigl(s_k(v)\partial_{z_k}+\bar s_k(v)\partial_{\bar z_k}\bigr).
\]
$A^v$ is the generator of the modes seen from a particle moving at constant
velocity $v$; its invariant law is $\pi$, the law of independent standard
complex Gaussians, for every $v$.

\paragraph{Initial data.} $(x_0,v_0)$ is deterministic and
\begin{equation}
  z_k(0)=\zeta_k+\eps\mu_k,\qquad
  (\zeta_k)\sim\pi,\quad \mu_k\in\mathbb{C}\ \text{deterministic},\quad
  \max_k|\mu_k|\le\mu_* .
  \label{eq:init}
\end{equation}
Two cases matter. \emph{Dressed start}: $\mu=0$, the Gibbs law of the field
given the particle (Lemma~\ref{lem:gibbs}(iii)); for $\kappa=0$ it is the
initial law of the companion note. \emph{Undressed start}: the field is in
its own equilibrium, $a_k(0)$ standard Gaussian in \eqref{eq:model}, which
is $\mu_k=\bar c_k/2T$ at $\kappa=1$; the particle then acquires its
polarization cloud during an initial transient of length $O(1/\gamma_k)$.
All results hold for both, with constants depending on $\mu_*$.

\paragraph{Constants.} Below, $C$ denotes a constant depending only on $d$,
$K$, $(c_k)$, $(\gamma_k)$, $(\omega_k)$, $T$ and $\mu_*$ (and on $q$ in
Section~\ref{sec:third}), and \emph{not} on $\eps\in(0,1]$,
$\kappa\in[0,1]$, $t$, $x_0$, $v_0$. Its value changes from line to line.
Write $|z|=\sum_k|z_k|$.

\begin{lemma}[Moments of the modes]\label{lem:mom}
Let $\hat a_k$ be the stationary Ornstein--Uhlenbeck process
$\dd\hat a_k=-(\gamma_k+i\omega_k)\hat a_k\dd t+\sqrt{2\gamma_k}\dd B_k$
with $\hat a_k(0)=e^{-ik\cdot x_0}\zeta_k$. Then, pathwise and for all
$t\ge0$,
\[
  |z_k(t)|\le|\hat a_k(t)|+\eps\rho_k,\qquad
  \rho_k=|\mu_k|+\frac{\kappa|c_k|}{2T}\Bigl(2+\frac{|\gamma_k+i\omega_k|}{\gamma_k}\Bigr).
\]
In particular \eqref{eq:kappa} has a unique global solution and, for every
$p\ge1$,
\[
  m_p:=\sup_{t\ge0}\E\bigl(1+|z(t)|\bigr)^p
  \le\E\Bigl(1+\sum_k(|\zeta_k|+\rho_k)\Bigr)^p<\infty .
\]
\end{lemma}

\begin{proof}
Let $a^\kappa_k=\tilde a_k-\eps\kappa\bar c_ke^{-ik\cdot x}/2T$. By
\eqref{eq:kappa},
$\dd a^\kappa_k=-(\gamma_k+i\omega_k)\tilde a_k\dd t+\sqrt{2\gamma_k}\dd B_k$,
so $\delta_k=a_k^\kappa-\hat a_k$ solves
$\dot\delta_k=-(\gamma_k+i\omega_k)\bigl(\delta_k
+\eps\kappa\bar c_ke^{-ik\cdot x}/2T\bigr)$ with
$|\delta_k(0)|\le\eps(|\mu_k|+\kappa|c_k|/2T)$. Hence
$|\delta_k(t)|\le\eps(|\mu_k|+\kappa|c_k|/2T)
+\eps\kappa|c_k||\gamma_k+i\omega_k|/(2T\gamma_k)$, and
$|z_k|=|\tilde a_k|\le|\hat a_k|+|\delta_k|+\eps\kappa|c_k|/2T$. Global
existence follows as in Lemma~\ref{lem:gibbs}(i). The moment bound uses
$\eps\le1$ and stationarity of $\hat a$.
\end{proof}

So, with back-reaction, the modes are no longer an Ornstein--Uhlenbeck
process independent of the orbit, but they differ from one by a
deterministically bounded amount of order $\eps$. This is the only a priori
information on the field that is used.

\section{Coefficients, the Poisson equation, and three identities}

Define
\begin{equation}
  D(v)=\frac12\sum_{k\in K}|c_k|^2\,\Re\!\Bigl(\frac1{\lambda_k(v)}\Bigr)\,k\otimes k,
  \qquad
  \Re\!\Bigl(\frac1{\lambda_k(v)}\Bigr)
   =\frac{\gamma_k}{\gamma_k^2+(\omega_k-k\cdot v)^2}.
  \label{eq:D}
\end{equation}
$D(v_0)$ is the matrix $D$ of the companion note for $f_k=-ic_kk$. Since
$\nabla_v(1/\lambda_k)=ik/\lambda_k^2$,
\begin{equation}
  \Div D(v)=\frac12\sum_k|c_k|^2|k|^2\,\Re\!\Bigl(\frac{i}{\lambda_k(v)^2}\Bigr)k
  =\sum_k|c_k|^2|k|^2\,
   \frac{\gamma_k\,(\omega_k-k\cdot v)}{\bigl(\gamma_k^2+(\omega_k-k\cdot v)^2\bigr)^2}\;k ,
  \label{eq:divD}
\end{equation}
where $(\Div D)_j=\sum_i\partial_{v_i}D_{ij}$. Define the drift
\begin{equation}
  m_\kappa(v)=\Div D(v)-\frac{\kappa}{T}\,D(v)\,v
  =\sum_k|c_k|^2\,k\,\Bigl[
    \frac{|k|^2\gamma_k(\omega_k-k\cdot v)}{(\gamma_k^2+(\omega_k-k\cdot v)^2)^2}
    -\frac{\kappa}{2T}\,\frac{\gamma_k\,(k\cdot v)}{\gamma_k^2+(\omega_k-k\cdot v)^2}\Bigr].
  \label{eq:m}
\end{equation}

\begin{lemma}\label{lem:coef}
The functions $1/\lambda_k(v)$ and $(k\cdot v)/\lambda_k(v)$ are bounded on
$\mathbb{R}^d$ with bounded derivatives of all orders. Consequently $D$,
$\Div D$, $v\mapsto D(v)v$ and $m_\kappa$ belong to
$C_b^\infty(\mathbb{R}^d)$, and $v\cdot D(v)v$ is bounded.
\end{lemma}

\begin{proof}
$|\lambda_k|\ge\gamma_k$ and
$|k\cdot v|\le|\omega_k|+|\omega_k-k\cdot v|\le|\omega_k|+|\lambda_k|$, so
$|k\cdot v|/|\lambda_k|\le1+|\omega_k|/\gamma_k$. Derivatives of
$\lambda_k^{-n}$ and of $(k\cdot v)\lambda_k^{-n}$ are linear combinations of
functions of the same two forms with $n\ge1$. Finally
$D(v)v=\frac12\sum|c_k|^2\Re(1/\lambda_k)(k\cdot v)k$ and
$v\cdot D(v)v=\frac12\sum|c_k|^2\Re(1/\lambda_k)(k\cdot v)^2
\le\frac12\sum|c_k|^2\gamma_k(k\cdot v)^2/|\lambda_k|^2$.
\end{proof}

\paragraph{Polynomial classes.} Let $\cP_n$ be the set of polynomials of
degree at most $n$ in $(z,\bar z)$ whose coefficients are functions in
$C_b^\infty(\mathbb{R}^d)$ of $v$. For $p\in\cP_n$ and every multi-index
$j$, $|\partial_v^jp(v,z)|\le C_{p,j}(1+|z|)^n$. Write $\E_\pi p$ for the
mean of $p(v,\cdot)$ under $\pi$, a function of $v$, and
\[
  \cK=G\cdot\nabla_v+\kappa S ,\qquad\text{so that}\qquad
  \cL=(v-v_0)\cdot\nabla_Y+A^v+\eps\,\cK .
\]
If $p\in\cP_n$ then $G\cdot\nabla_vp\in\cP_{n+1}$, and $Sp$ is a polynomial
of degree $n-1$ whose coefficients are $(k\cdot v)$ times functions in
$C_b^\infty$; hence
\begin{equation}
  |\cK p(v,z)|\le C_p\,(1+|v|)\,(1+|z|)^{n+1}.
  \label{eq:Kbound}
\end{equation}

\begin{lemma}[Poisson equation]\label{lem:poisson}
Let $p\in\cP_n$ with $\E_\pi p=0$ for every $v$. There is a unique
$Qp\in\cP_n$ with $\E_\pi Qp=0$ and $A^vQp=-p$. The map $Q$ is linear, acts
on the coefficients of $p$ by multiplication with functions in $C_b^\infty$
determined by $(\gamma_k,\omega_k,k)$, and preserves parity in $z$.
\end{lemma}

\begin{proof}
For multi-indices $a,b\in\mathbb{N}^K$,
\[
  A^v(z^a\bar z^b)=-\Lambda_{ab}(v)\,z^a\bar z^b
   +\sum_k2\gamma_ka_kb_k\,z^{a-e_k}\bar z^{b-e_k},
  \qquad
  \Lambda_{ab}=\sum_k\bigl(a_k\lambda_k+b_k\bar\lambda_k\bigr).
\]
For $(a,b)\ne0$, $\Re\Lambda_{ab}\ge\min_k\gamma_k>0$ and
$\nabla_v\Lambda_{ab}$ is constant, so $1/\Lambda_{ab}\in C_b^\infty$. By
induction on $n$ we claim: for every $p\in\cP_n$ there is $u\in\cP_n$
without constant term such that $A^vu+p$ does not depend on $z$. For $n=0$
take $u=0$. For $n\ge1$ let $\sum_{|a|+|b|=n}p_{ab}z^a\bar z^b$ be the top
degree part of $p$ and $u_n=\sum p_{ab}\Lambda_{ab}^{-1}z^a\bar z^b$; then
$A^vu_n+p\in\cP_{n-1}$ and the induction hypothesis applies to it. Next,
$\int A^v(z^a\bar z^b)\dd\pi=0$: since
$\int z^a\bar z^b\dd\pi=\delta_{ab}\,a!$, both terms vanish if $a\ne b$, and
if $a=b$ the sum is $-2\sum_k\gamma_ka_k\,a!+\sum_k2\gamma_ka_k^2(a-e_k)!=0$.
Therefore $A^vu+p=\E_\pi p=0$, and $Qp=u-\E_\pi u$. If $u\in\cP_n$ and
$A^vu=0$, the top degree part of $u$ vanishes unless it is constant, so $u$
is a function of $v$; this gives uniqueness. Linearity, the form of the
coefficients, and parity are clear from the construction.
\end{proof}

Define the vector-valued first corrector and its companion
\begin{equation}
  \Psi(v,z)=\sum_k\Re\Bigl(\frac{f_k\,z_k}{\lambda_k(v)}\Bigr),\qquad
  \Psi'(v,z)=\sum_k\Re\Bigl(\frac{f_k\,z_k}{\lambda_k(v)^2}\Bigr).
  \label{eq:Psi}
\end{equation}
Compared with the companion note, the resonance denominator is evaluated at
the \emph{current} velocity $v$, not at $v_0$. This removes the need for the
corrector $\alpha_2$ of that note and is what makes the kinetic time
accessible.

\begin{lemma}[Identities]\label{lem:id}
For all $v$ and $z$:
\begin{enumerate}
  \item[(a)] $A^v\Psi=-G$ and $A^v\Psi'=-\Psi$;
  \item[(b)] $S\Psi=-\dfrac1T\,D(v)\,v$ \quad(polarization drag);
  \item[(c)] $\E_\pi[G_i\Psi_j]=D_{ij}(v)$;
  \item[(d)] $\E_\pi\bigl[(G\cdot\nabla_v)\Psi\bigr]=\Div D(v)$;
  \item[(e)] $|\Psi|+|\Psi'|+|v\cdot\Psi|\le C|z|$, \
    $|S\Psi'|\le C$, and
    $\bigl|\sum_jw_j\,(G\cdot\nabla_v)\Psi_j\bigr|\le C(1+|v_0|)|z|^2$ for
    $w=v-v_0$.
\end{enumerate}
\end{lemma}

\begin{proof}
(a) $A^vz_k=-\lambda_k(v)z_k$. (b) Since
$\Psi=\frac12\sum_k(f_kz_k/\lambda_k+\text{c.c.})$,
\[
  S\Psi=\sum_k\Re\Bigl(\frac{f_ks_k}{\lambda_k}\Bigr)
  =\sum_k\Re\Bigl(\frac{(-ic_kk)\,(-i)(k\cdot v)\bar c_k}{2T\lambda_k}\Bigr)
  =-\frac1{2T}\sum_k|c_k|^2\Re\Bigl(\frac1{\lambda_k}\Bigr)k\,(k\cdot v).
\]
(c) Under $\pi$, $\E z_l\bar z_k=\delta_{lk}$ and $\E z_lz_k=0$, so
$\E_\pi[G_i\Psi_j]=\frac14\sum_k\bigl(f_{k,i}\bar f_{k,j}/\bar\lambda_k
+\bar f_{k,i}f_{k,j}/\lambda_k\bigr)
=\frac12\sum_k|c_k|^2k_ik_j\Re(1/\lambda_k)$.
(d) $\partial_{v_i}\Psi_j=\sum_k\Re(ik_if_{k,j}z_k/\lambda_k^2)$, hence
$\E_\pi[G_i\partial_{v_i}\Psi_j]
=\frac12\sum_k\Re\bigl(i(k\cdot\bar f_k)f_{k,j}/\lambda_k^2\bigr)$, and
$i(k\cdot\bar f_k)f_{k,j}=i\,(i\bar c_k|k|^2)(-ic_kk_j)=i|c_k|^2|k|^2k_j$;
compare \eqref{eq:divD}.
(e) $v\cdot f_k/\lambda_k=-ic_k(k\cdot v)/\lambda_k$ and
$f_ks_k/\lambda_k^2$ are bounded by Lemma~\ref{lem:coef}. Further
$\sum_jw_j(G\cdot\nabla_v)\Psi_j
=\sum_k\Re\bigl(i(G\cdot k)(w\cdot f_k)z_k/\lambda_k^2\bigr)$ and
$|w\cdot f_k|/|\lambda_k|^2\le|c_k|(|k\cdot v|+|k||v_0|)/|\lambda_k|^2$.
\end{proof}

Identities (b) and (d) are exact and hold for every $v$; they are the
source of the two parts of the drift $m_\kappa$. They were also verified
symbolically ($d=1$, two modes, exact rational arithmetic at randomly
chosen parameter values), together with Lemmas~\ref{lem:rem},
\ref{lem:diff} and the expansions in the proofs of
Theorems~\ref{thm:drift} and~\ref{thm:kin}. The symbolic checks were
scratch computations and are not part of the repository.

\begin{remark}[Which drift without back-reaction]
For $\kappa=0$ the drift is $\Div D$, so the limiting generator is
$D:\nabla^2+\Div D\cdot\nabla=\nabla\cdot(D\nabla)$, in divergence form, and
not $D:\nabla^2$. For a general (non-potential) field $f_k$ the same
computation gives the drift
$\frac12\sum_k\Re\bigl(i(k\cdot\bar f_k)f_k/\lambda_k^2\bigr)$, which need
not equal $\Div D$; the divergence form is tied to the field being a
gradient. Only potential fields are treated in this note.
\end{remark}

\section{A priori bounds}

\begin{lemma}[Dynkin's formula]\label{lem:dynkin}
Let $F(\tau,Y,v,z)$ be continuous on $[0,t]\times\mathbb{R}^d\times
\mathbb{R}^d\times\mathbb{C}^K$, once continuously differentiable in
$(\tau,Y,v)$ and twice in $z$, with $|F|$, $|\partial_zF|$ and
$|(\partial_\tau+\cL)F|$ bounded by $C_F(1+|Y|+|v|)^a(1+|z|)^b$ for some
$a,b$. Then
\[
  \E F(t,Y_t,v_t,z_t)-\E F(0,0,v_0,z_0)
  =\E\int_0^t(\partial_\tau+\cL)F(\tau,Y_\tau,v_\tau,z_\tau)\,\dd\tau .
\]
\end{lemma}

\begin{proof}
$Y$ and $v$ have finite variation, so It\^o's formula applies and gives the
identity up to a stochastic integral against the $W_k$ with integrands
$\sqrt{2\gamma_k}\,\partial_{z_k}F$ and conjugates. Since
$|v_\tau-v_0|\le\eps\int_0^\tau|G(z_\sigma)|\dd\sigma$ and
$|Y_\tau|\le\tau\sup_{\sigma\le\tau}|v_\sigma-v_0|$, Minkowski's inequality
and Lemma~\ref{lem:mom} give $\sup_{\tau\le t}\E(1+|Y_\tau|+|v_\tau|)^p<\infty$
for every $p$. By H\"older's inequality the integrands are square integrable
on $[0,t]$, the stochastic integral is a martingale, and all terms are
integrable.
\end{proof}

\begin{lemma}[Diffusive bound]\label{lem:diff}
For all $t\ge0$,
\[
  \E|v_t-v_0|^2\le C\,(1+|v_0|)\,\eps^2\,(1+t).
\]
\end{lemma}

\begin{proof}
Let $w=v-v_0$ and $h=|w|^2+2\eps\,w\cdot\Psi(v,z)$. Using
Lemma~\ref{lem:id}(a),(b),
\[
  \cL h=2\eps\,G\cdot w+2\eps^2\bigl[G\cdot\Psi+\textstyle\sum_jw_j(G\cdot\nabla_v)\Psi_j\bigr]
     -2\eps\,w\cdot G-\frac{2\eps^2\kappa}{T}\,w\cdot D(v)v .
\]
The terms of order $\eps$ cancel. By Lemma~\ref{lem:id}(e) and
Lemma~\ref{lem:coef} the rest is bounded in absolute value by
$C\eps^2(1+|v_0|)(1+|z|)^2$. Lemma~\ref{lem:dynkin} with $h(0)=0$ gives
$\E|w_t|^2\le2\eps\,\E|w_t||\Psi|+C\eps^2(1+|v_0|)\,m_2\,t$, and
$2\eps|w||\Psi|\le\frac12|w|^2+2\eps^2|\Psi|^2$ with $\E|w_t|^2<\infty$.
\end{proof}

This is the diffusive a priori bound that the companion note did not have
(its proof uses only $|w_\tau|\lesssim\eps\tau$, as the remark closing that proof points out). It holds
for all times, including the kinetic time, because the corrector $\Psi$ is
exact in $v$ and $w\cdot\Psi$ is controlled by the resonance denominators.

\begin{lemma}[Equal-time correlations]\label{lem:equal}
Let $\phi\in C_b^1(\mathbb{R}^d;\mathbb{C})$ be such that
$(k\cdot v)\phi(v)$ is bounded. Then for all $t\ge0$
\[
  \bigl|\E[\phi(v_t)\,z_k(t)]\bigr|\le C_\phi\,\eps,
\]
with $C_\phi$ depending on the data listed above and on
$\nrm{\phi}{\infty}$, $\nrm{\nabla\phi}{\infty}$,
$\nrm{(k\cdot v)\phi}{\infty}$. In particular
$|\E\Psi(v_t,z_t)|\le C\eps$.
\end{lemma}

\begin{proof}
Fix $t$ and let
$h(\tau,v,z)=\phi(v)\,e^{-\lambda_k(v)(t-\tau)}z_k$. Since
$\nabla_ve^{-\lambda_k(v)(t-\tau)}=ik(t-\tau)e^{-\lambda_k(v)(t-\tau)}$ and
$A^vz_k=-\lambda_kz_k$,
\[
  (\partial_\tau+\cL)h=\eps\,e^{-\lambda_k(v)(t-\tau)}
  \Bigl[G\cdot\bigl(\nabla\phi+i(t-\tau)k\,\phi\bigr)z_k+\kappa\,\phi\,s_k\Bigr].
\]
By Lemma~\ref{lem:dynkin} and $|e^{-\lambda_k(t-\tau)}|=e^{-\gamma_k(t-\tau)}$,
\[
  \bigl|\E[\phi(v_t)z_k(t)]-e^{-\lambda_k(v_0)t}\phi(v_0)\,\eps\mu_k\bigr|
  \le\eps\int_0^te^{-\gamma_k(t-\tau)}\bigl[(\nrm{\nabla\phi}{\infty}
   +|k|(t-\tau)\nrm{\phi}{\infty})\,C\,m_2+\nrm{\phi s_k}{\infty}\bigr]\dd\tau ,
\]
which is at most $C_\phi\eps$ uniformly in $t$. For $\Psi$ take
$\phi=f_{k,j}/\lambda_k$ and its conjugate.
\end{proof}

\section{The friction: mean drift and the Einstein relation}

\begin{theorem}[Mean drift]\label{thm:drift}
For all $\eps\in(0,1]$, $\kappa\in[0,1]$, $t\ge0$ and initial data
\eqref{eq:init},
\begin{align}
  \bigl|\E v_t-v_0-\eps^2t\,m_\kappa(v_0)\bigr|
   &\le C(1+|v_0|)\bigl[\eps^2+\eps^3(t+t^{3/2})\bigr],
   \label{eq:drift}\\
  \bigl|\E Y_t-\tfrac12\eps^2t^2\,m_\kappa(v_0)\bigr|
   &\le C(1+|v_0|)\bigl[\eps^2t+\eps^3(t^2+t^{5/2})\bigr],
   \label{eq:driftY}
\end{align}
with $m_\kappa$ given by \eqref{eq:m}. The estimate \eqref{eq:drift}
identifies the drift for $1\ll t\ll\eps^{-2}$. With $\eps=N^{-1/2}$ and
$t=sN^{1/3}$,
\[
  \E Y_t=\tfrac12\,s^2\,m_\kappa(v_0)\,N^{-1/3}+O\bigl((1+|v_0|)(s+s^{5/2})N^{-2/3}\bigr).
\]
\end{theorem}

The drift has two parts. $\Div D(v_0)$ is present in the prescribed field
($\kappa=0$). The back-reaction adds exactly $-D(v_0)v_0/T$.

\begin{proof}
Let $\Phi=(G\cdot\nabla_v)\Psi-\Div D\in\cP_2$ (componentwise). By
Lemma~\ref{lem:id}(d), $\E_\pi\Phi=0$; let $\Xi=Q\Phi\in\cP_2$
(Lemma~\ref{lem:poisson}). Using $\cL=A^v+\eps\cK$ on functions of $(v,z)$
and Lemma~\ref{lem:id}(a),(b),
\begin{align*}
  \cL v&=\eps G,\qquad
  \cL(\eps\Psi)=-\eps G+\eps^2\bigl[(G\cdot\nabla_v)\Psi-\kappa T^{-1}D(v)v\bigr]
   =-\eps G+\eps^2\bigl[m_\kappa(v)+\Phi\bigr],\\
  \cL(\eps^2\Xi)&=-\eps^2\Phi+\eps^3\cK\Xi .
\end{align*}
Hence the exact identity
\begin{equation}
  \cL\bigl(v+\eps\Psi+\eps^2\Xi\bigr)=\eps^2m_\kappa(v)+\eps^3\cK\Xi .
  \label{eq:exactdrift}
\end{equation}
By Lemma~\ref{lem:dynkin},
\[
  \E v_t-v_0=\eps^2\!\int_0^t\!\E\,m_\kappa(v_\tau)\dd\tau
   -\eps\,\E\Psi(v_t,z_t)+\eps\,\E\Psi(v_0,z_0)
   -\eps^2\,\E\bigl[\Xi(v_t,z_t)-\Xi(v_0,z_0)\bigr]
   +\eps^3\!\int_0^t\!\E\,\cK\Xi\,\dd\tau .
\]
We bound each term. By Lemma~\ref{lem:equal}, $|\eps\E\Psi(v_t,z_t)|\le
C\eps^2$. Since $\Psi$ is linear in $z$,
$\eps\E\Psi(v_0,z_0)=\eps^2\sum_k\Re(f_k\mu_k/\lambda_k(v_0))$, of modulus
at most $C\eps^2$ (and zero for the dressed start). $\eps^2\E|\Xi|\le
C\eps^2m_2$. By \eqref{eq:Kbound}, the Cauchy--Schwarz inequality and
Lemma~\ref{lem:diff},
\[
  \E|\cK\Xi|\le C\bigl(\E(1+|v_\tau|)^2\bigr)^{1/2}m_6^{1/2}
  \le C\bigl(1+|v_0|+(1+|v_0|)^{1/2}\eps(1+\tau)^{1/2}\bigr),
\]
so the last term is at most $C(1+|v_0|)\eps^3(t+t^{3/2})$. Finally
$m_\kappa$ is Lipschitz (Lemma~\ref{lem:coef}), so by Lemma~\ref{lem:diff}
\[
  \Bigl|\int_0^t\E\,m_\kappa(v_\tau)\dd\tau-t\,m_\kappa(v_0)\Bigr|
  \le C\int_0^t\bigl(\E|v_\tau-v_0|^2\bigr)^{1/2}\dd\tau
  \le C(1+|v_0|)^{1/2}\eps\,(t+t^{3/2}).
\]
This proves \eqref{eq:drift}; \eqref{eq:driftY} follows by integrating in
time, since $\E Y_t=\int_0^t(\E v_\tau-v_0)\dd\tau$.
\end{proof}

\begin{remark}[Dressing transient]
The term $\eps\E\Psi(v_0,z_0)=\eps^2\sum_k\Re(f_k\mu_k/\lambda_k(v_0))$ is
the only place where the two initial conditions differ in the proof. For the
undressed start at $\kappa=1$ it equals
$-\frac{\eps^2}{2T}\sum_k|c_k|^2k\,
(\omega_k-k\cdot v_0)/(\gamma_k^2+(\omega_k-k\cdot v_0)^2)$, a velocity kick
of order $\eps^2$ acquired while the polarization cloud forms. It is inside
the error term $C\eps^2$ of \eqref{eq:drift}; we do not claim that it is the
whole $O(\eps^2)$ difference between the two starts.
\end{remark}

\begin{proposition}[Einstein relation]\label{prop:einstein}
Let $\cL_\kappa=D:\nabla_v^2+m_\kappa\cdot\nabla_v$ and
$M_\Theta(v)=(2\pi\Theta)^{-d/2}e^{-|v|^2/2\Theta}$. Then
\begin{enumerate}
  \item[(i)] $\cL_\kappa\varphi=\nabla\cdot(D\nabla\varphi)
     -\kappa T^{-1}(Dv)\cdot\nabla\varphi$; the back-reaction drift is
     $-T^{-1}Dv=D\,\nabla\log M_T$.
  \item[(ii)] The formal adjoint is
     $\cL_\kappa^*f=\nabla\cdot\bigl(D(\nabla f+\kappa T^{-1}vf)\bigr)$.
  \item[(iii)] For $\kappa=1$,
     $\cL_1\varphi=M_T^{-1}\nabla\cdot(M_TD\nabla\varphi)$, and for
     $\varphi,\psi\in C_b^2(\mathbb{R}^d)$,
     \[
       \int\psi\,\cL_1\varphi\,M_T\,\dd v=-\int\nabla\psi\cdot D\nabla\varphi\,M_T\,\dd v .
     \]
     So $\cL_1$ is symmetric and nonpositive in $L^2(M_T)$,
     $\cL_1^*M_T=0$, and $\int\cL_1\varphi\,M_T=0$.
  \item[(iv)] For $\kappa\in(0,1]$, $\cL_\kappa^*M_{T/\kappa}=0$. For
     $\kappa=0$, $\cL_0=\nabla\cdot(D\nabla)$ is symmetric with respect to
     Lebesgue measure and $\cL_0^*M_\Theta=-\Theta^{-1}\nabla\cdot(DvM_\Theta)$.
\end{enumerate}
\end{proposition}

\begin{proof}
(i) is the definition of $m_\kappa$ and
$\nabla\cdot(D\nabla\varphi)=D:\nabla^2\varphi+\Div D\cdot\nabla\varphi$
($D$ is symmetric). (ii) Integrate by parts. (iii)
$M_T^{-1}\nabla\cdot(M_TD\nabla\varphi)=\nabla\cdot(D\nabla\varphi)
+\nabla\log M_T\cdot D\nabla\varphi$ and $\nabla\log M_T=-v/T$. The
integration by parts has no boundary term because $D$ and the derivatives of
$\varphi,\psi$ are bounded and $M_T$ is Gaussian. (iv)
$\nabla M_\Theta=-vM_\Theta/\Theta$, so
$\cL_\kappa^*M_\Theta=(\kappa/T-1/\Theta)\nabla\cdot(DvM_\Theta)$.
\end{proof}

Statement (iii) is the fluctuation--dissipation relation of the model: the
friction produced by the back-reaction (Lemma~\ref{lem:id}(b)) and the
diffusion produced by the field fluctuations (Lemma~\ref{lem:id}(c)) involve
the same matrix $D(v)$, with ratio $1/T$, at every velocity and not only on
average. The stationary velocity law $M_T$ is the velocity marginal of the
Gibbs measure of Lemma~\ref{lem:gibbs}. Without back-reaction no Maxwellian
is stationary: the particle is heated by the field. (Item (iv) for
$0<\kappa<1$ is an algebraic observation about the interpolating family.)

\section{Back-reaction is negligible at \texorpdfstring{$N^{1/3}$}{N\textasciicircum(1/3)}}
\label{sec:third}

\begin{theorem}\label{thm:third}
Let $q\in\mathbb{R}^d$, $\eps\in(0,1]$, $\kappa\in[0,1]$, and $t\ge0$ with
$\eps^2t\le1$. For initial data \eqref{eq:init},
\begin{equation}
  \Bigl|\E\,e^{iq\cdot Y_t}-\exp\!\Bigl(-\tfrac13\eps^2t^3\,q\cdot D(v_0)q\Bigr)\Bigr|
  \le C(1+|v_0|)\Bigl[\eps+\eps^2(t+t^2)+\eps^3(t^3+t^{7/2})
     +\eps^4(t^4+t^5)+\eps^5t^6\Bigr],
  \label{eq:third}
\end{equation}
with $D$ given by \eqref{eq:D}, the same matrix as without back-reaction,
and $C$ independent of $\kappa$. In particular, with $\eps=N^{-1/2}$ and
$t=sN^{1/3}$,
\[
  \E\,e^{iq\cdot Y_t}=\exp\!\Bigl(-\tfrac13s^3\,q\cdot D(v_0)q\Bigr)
   +O\bigl((1+|v_0|)(1+s^6)\,N^{-1/3}\bigr),
\]
and the right-hand side of \eqref{eq:third} tends to zero as $N\to\infty$
whenever $t\ll N^{2/5}$.
\end{theorem}

\begin{corollary}\label{cor:law}
For the coupled model ($\kappa=1$) with either the dressed or the undressed
start, as $N\to\infty$ with $t=sN^{1/3}$, $Y_t$ converges in law to a
centred Gaussian vector with covariance $\tfrac23s^3D(v_0)$: the limit is
the same as in the prescribed field.
\end{corollary}

\begin{corollary}\label{cor:rate}
For the prescribed-field model of the companion note with a potential field
$f_k=-ic_kk$ (this is $\kappa=0$, $\mu=0$), the error at $t=sN^{1/3}$ is
$O(N^{-1/3})$ instead of $O(N^{-1/6})$, and the range of validity is
$t\ll N^{2/5}$ instead of $t\ll N^{3/8}$.
\end{corollary}

\subsection{Proof of Theorem~\ref{thm:third}}

Fix $q$ and $t$; write $\theta=t-\tau$, $w=v-v_0$, $g=q\cdot G$,
$d(v)=q\cdot D(v)q$, $d_0=d(v_0)$. Let
\[
  U(\tau,Y,v)=\exp\!\Bigl(iq\cdot Y+i\theta\,q\cdot w-\tfrac13\eps^2\theta^3d_0\Bigr),
\]
so that $|U|\le1$,
$(\partial_\tau+w\cdot\nabla_Y)U=\eps^2\theta^2d_0\,U$ and
$\nabla_vU=i\theta q\,U$. Consequently, for a function $p(\tau,v,z)$,
\begin{equation}
  (\partial_\tau+\cL)(pU)
  =U\bigl[\partial_\tau p+\eps^2\theta^2d_0\,p+i\eps\theta\,g\,p
     +\eps\,\cK p+A^vp\bigr].
  \label{eq:pU}
\end{equation}
\paragraph{Correctors.} Let
\[
  \alpha=q\cdot\Psi,\qquad\alpha'=q\cdot\Psi',\qquad
  \beta=-Q(\alpha g-d),\qquad c=-Q(\beta g).
\]
Here $\alpha,\alpha'\in\cP_1$ with $A^v\alpha=-g$, $A^v\alpha'=-\alpha$
(Lemma~\ref{lem:id}(a)); $\E_\pi[\alpha g]=d(v)$ by Lemma~\ref{lem:id}(c),
so $\beta\in\cP_2$ is well defined, even in $z$, with $A^v\beta=\alpha g-d$;
and $\beta g$ is odd in $z$, so $\E_\pi[\beta g]=0$ and $c\in\cP_3$ is well
defined, odd, with $A^vc=\beta g$. Explicitly, with
$A_k=q\cdot f_k/\lambda_k(v)$ and $B_k=q\cdot f_k$,
\[
  \beta=-\frac14\sum_{k,l}\Bigl(\frac{A_kB_l}{\lambda_k+\lambda_l}\,z_kz_l
   +\frac{A_k\bar B_l}{\lambda_k+\bar\lambda_l}\,(z_k\bar z_l-\delta_{kl})
   +\text{c.c.}\Bigr),
\]
the formula of the companion note with $\lambda_k=\lambda_k(v)$. Set
\begin{equation}
  U^\eps=U\Bigl[1+i\eps\theta\,\alpha+\eps^2\theta^2\beta
    -i\eps\,\alpha'+i\eps^3\theta^3d_0\,\alpha'-i\eps^3\theta^3c\Bigr].
\end{equation}
The first two correctors are those of the companion note (with $v$ in place
of $v_0$). The last three remove the terms that limited its rate to
$N^{-1/6}$.

\begin{lemma}\label{lem:rem}
$(\partial_\tau+\cL)U^\eps=UR$ with
\begin{align*}
  R={}&\eps^2\theta^2\bigl(d_0-d(v)\bigr)
     +i\eps^2\theta\,\cK\alpha-2\eps^2\theta\,\beta+\eps^2\theta\,g\alpha'
     -i\eps^2\,\cK\alpha' \\
   &+\eps^3\theta^2\,\cK\beta-4i\eps^3\theta^2d_0\,\alpha'+3i\eps^3\theta^2c \\
   &+\eps^4\theta^4d_0\,\beta-\eps^4\theta^4d_0\,g\alpha'+\eps^4\theta^4gc
     +i\eps^4\theta^3d_0\,\cK\alpha'-i\eps^4\theta^3\,\cK c \\
   &+i\eps^5\theta^5d_0^2\,\alpha'-i\eps^5\theta^5d_0\,c .
\end{align*}
\end{lemma}

\begin{proof}
Apply \eqref{eq:pU} to each term, with $\partial_\tau\theta=-1$. The
bracket in \eqref{eq:pU} is, for
\begin{align*}
  p&=1: &&\eps^2\theta^2d_0+i\eps\theta g;\\
  p&=i\eps\theta\alpha: &&-i\eps\alpha+i\eps^3\theta^3d_0\alpha
     -\eps^2\theta^2g\alpha+i\eps^2\theta\,\cK\alpha-i\eps\theta g;\\
  p&=\eps^2\theta^2\beta: &&-2\eps^2\theta\beta+\eps^4\theta^4d_0\beta
     +i\eps^3\theta^3g\beta+\eps^3\theta^2\cK\beta+\eps^2\theta^2(\alpha g-d(v));\\
  p&=-i\eps\alpha': &&-i\eps^3\theta^2d_0\alpha'+\eps^2\theta g\alpha'
     -i\eps^2\cK\alpha'+i\eps\alpha;\\
  p&=i\eps^3\theta^3d_0\alpha': &&-3i\eps^3\theta^2d_0\alpha'
     +i\eps^5\theta^5d_0^2\alpha'-\eps^4\theta^4d_0g\alpha'
     +i\eps^4\theta^3d_0\cK\alpha'-i\eps^3\theta^3d_0\alpha;\\
  p&=-i\eps^3\theta^3c: &&3i\eps^3\theta^2c-i\eps^5\theta^5d_0c
     +\eps^4\theta^4gc-i\eps^4\theta^3\cK c-i\eps^3\theta^3\beta g.
\end{align*}
On summing, $\pm i\eps\theta g$, $\pm i\eps\alpha$,
$\pm i\eps^3\theta^3d_0\alpha$, $\pm i\eps^3\theta^3\beta g$ and
$\pm\eps^2\theta^2\alpha g$ cancel, and
$\eps^2\theta^2d_0-\eps^2\theta^2d(v)$ remains from the first and third
lines. The other terms are those listed.
\end{proof}

\begin{proof}[Proof of Theorem~\ref{thm:third}]
$U^\eps$ is smooth and satisfies the growth conditions of
Lemma~\ref{lem:dynkin}, so
$\E U^\eps(t)-\E U^\eps(0)=\E\int_0^tUR\,\dd\tau$.

\emph{Endpoint $\tau=t$.} $\theta=0$ and
$U^\eps=e^{iq\cdot Y_t}(1-i\eps\alpha'(v_t,z_t))$, so
$|\E U^\eps(t)-\E e^{iq\cdot Y_t}|\le\eps\E|\alpha'|\le C\eps$.

\emph{Endpoint $\tau=0$.} $U^\eps(0)=e^{-\eps^2t^3d_0/3}[1+P(v_0,z_0)]$ with
$P=i\eps t\alpha+\eps^2t^2\beta-i\eps\alpha'+i\eps^3t^3d_0\alpha'-i\eps^3t^3c$.
Each of $\alpha,\alpha',\beta,c$ has zero mean under $\pi$, and for a
polynomial $p$ with $\E_\pi p=0$,
$|\E p(v_0,\zeta+\eps\mu)|\le C\eps$ (every term of
$p(\zeta+\eps\mu)-p(\zeta)$ carries a factor $\eps$). Hence
$|\E U^\eps(0)-e^{-\eps^2t^3d_0/3}|\le C(\eps^2+\eps^2t+\eps^3t^2+\eps^4t^3)$,
and it vanishes for the dressed start.

\emph{The remainder.} $|U|\le1$. By Lemma~\ref{lem:diff} and $\eps^2t\le1$,
$\E(1+|v_\tau|)^2\le C(1+|v_0|)^2$ for $\tau\le t$. For $p\in\cP_n$,
$\E|p|\le Cm_n$, and by \eqref{eq:Kbound} and Cauchy--Schwarz
$\E|\cK p|\le C(1+|v_0|)$. Since $d$ is Lipschitz (Lemma~\ref{lem:coef}),
$\E|d_0-d(v_\tau)|\le C(1+|v_0|)^{1/2}\eps(1+\tau)^{1/2}$ by
Lemma~\ref{lem:diff}. With
$\int_0^t\theta^a\dd\tau=t^{a+1}/(a+1)$ and
$\int_0^t\theta^2(1+\tau)^{1/2}\dd\tau\le t^3+t^{7/2}$, the time integrals of
the five groups of terms of $R$ are bounded by $C(1+|v_0|)$ times
\begin{center}
\renewcommand{\arraystretch}{1.3}
\begin{tabular}{ll}
  terms of $R$ & bound \\ \hline
  $\eps^2\theta^2(d_0-d(v))$ & $\eps^3(t^3+t^{7/2})$ \\
  $\eps^2\theta\,(\cK\alpha,\ \beta,\ g\alpha')$ and $\eps^2\cK\alpha'$
     & $\eps^2(t^2+t)$ \\
  $\eps^3\theta^2\,(\cK\beta,\ d_0\alpha',\ c)$ & $\eps^3t^3$ \\
  $\eps^4\theta^4\,(d_0\beta,\ d_0g\alpha',\ gc)$ and
  $\eps^4\theta^3\,(d_0\cK\alpha',\ \cK c)$ & $\eps^4(t^5+t^4)$ \\
  $\eps^5\theta^5\,(d_0^2\alpha',\ d_0c)$ & $\eps^5t^6$ \\
\end{tabular}
\end{center}
Adding the endpoint terms gives \eqref{eq:third}. At $t=s\eps^{-2/3}$:
$\eps^2t^2=s^2\eps^{2/3}$, $\eps^3t^{7/2}=s^{7/2}\eps^{2/3}$,
$\eps^4t^5=s^5\eps^{2/3}$, $\eps^5t^6=s^6\eps$, and the remaining terms are
smaller. If $\eps\to0$ and $\eps^4t^5\to0$, that is
$t\ll\eps^{-4/5}=N^{2/5}$, then every term tends to zero
($\eps^3t^{7/2}=(\eps^4t^5)^{7/10}\eps^{1/5}$ and
$\eps^5t^6=(\eps^4t^5)^{6/5}\eps^{1/5}$).
\end{proof}

\begin{remark}[Where the back-reaction enters]\label{rem:where}
$\kappa$ appears in $R$ only through $\kappa S$ in $\cK$. The largest such
term is, by Lemma~\ref{lem:id}(b),
\[
  i\eps^2\theta\,\kappa\,S\alpha=-\,i\eps^2\theta\,\frac{\kappa}{T}\,q\cdot D(v)\,v ,
\]
which does not depend on $z$ and whose time integral is at most
$\kappa\,\eps^2t^2\sup_v|q\cdot D(v)v|/2T$. This is the friction acting on
the phase: a displacement $-\kappa\eps^2t^2D(v_0)v_0/2T$, of order
$N^{-1/3}$ at $t\sim N^{1/3}$, in agreement with \eqref{eq:driftY} and with
the estimate $t^2/N$ in gap~2 of the companion note. All other
$\kappa$-dependent terms are $O(\eps^2t+\eps^3t^3+\eps^4t^4)$. So
back-reaction changes the characteristic function at order $N^{-1/3}$ and
not before; it does not change $D$.
\end{remark}

\begin{remark}[The next order is not identified]\label{rem:next}
One cannot read the $N^{-1/3}$ coefficient of $\E e^{iq\cdot Y_t}$ off
\eqref{eq:driftY}. In $R$ there are other terms of the same order at
$t=s\eps^{-2/3}$: $\eps^2\theta\,g\alpha'$ has non-zero $\pi$-mean (it is
the $O(t^2)$ correction to the variance in the proposition on the first
iterate in the companion note); $\eps^2\theta^2(d_0-d(v))$ and $\eps^4\theta^4gc$ contribute at order
$\eps^4t^5$ through the dependence of $D$ on the diffusing velocity. The
Monte Carlo run below shows an imaginary part for $\kappa=0$ that is not
equal to $e^{-s^3d_0/3}\,\tfrac12\eps^2t^2\,q\cdot\Div D$, consistent with
this. Identifying the full $N^{-1/3}$ term would need one more layer of
correctors, polynomial in $w$; this was not done.
\end{remark}

\section{The kinetic time}

Let $\Sigma(v)$ be the $d\times|K|$ matrix with columns
$\bigl(\tfrac12\Re(1/\lambda_k(v))\bigr)^{1/2}|c_k|\,k$, so that
$\Sigma\Sigma^{\!\top}=D$ and $\Sigma\in C_b^\infty$ (because
$\Re(1/\lambda_k)=\gamma_k/|\lambda_k|^2>0$ is smooth with a smooth square
root $\gamma_k^{1/2}/|\lambda_k|$). The diffusion
\begin{equation}
  \dd V=m_\kappa(V)\,\dd\tau+\sqrt2\,\Sigma(V)\,\dd\mathcal{W}
  \label{eq:FPsde}
\end{equation}
has generator $\cL_\kappa$ of Proposition~\ref{prop:einstein}.

\begin{theorem}[Kinetic limit with friction]\label{thm:kin}
Let $\tau_*>0$ and let $u:[0,\tau_*]\times\mathbb{R}^d\to\mathbb{R}$ be
continuous, differentiable in $\tau$, with $\partial_v^ju$ continuous and
bounded on $[0,\tau_*]\times\mathbb{R}^d$ for $|j|\le4$, and
\[
  \partial_\tau u=\cL_\kappa u
  =\nabla\cdot(D\nabla u)-\frac{\kappa}{T}(Dv)\cdot\nabla u,
  \qquad u(0,\cdot)=\varphi .
\]
Let $\mathcal{N}=\sup_{\tau\le\tau_*}\sum_{|j|\le4}\nrm{\partial_v^ju(\tau)}{\infty}$.
Then for all $\eps\in(0,1]$, $\kappa\in[0,1]$, initial data \eqref{eq:init},
and $t\ge0$ with $\tau=\eps^2t\le\tau_*$,
\begin{equation}
  \bigl|\E\varphi(v_t)-u(\eps^2t,v_0)\bigr|
  \le C\,\mathcal{N}\,(1+|v_0|)\,(1+\tau)^{3/2}\,\eps .
\end{equation}
With $\eps=N^{-1/2}$ this is an error $O(N^{-1/2})$ at the kinetic time
$t=\tau N$.
\end{theorem}

\begin{remark}[Existence and regularity of $u$]\label{rem:reg}
The coefficients $m_\kappa$ and $\Sigma$ of \eqref{eq:FPsde} are in
$C_b^\infty$ (Lemma~\ref{lem:coef}). For such coefficients it is classical
that, for $\varphi\in C_b^4$, $u(\tau,v)=\E\varphi(V_\tau^v)$ has four
continuous bounded $v$-derivatives, locally uniformly in $\tau$, and solves
the backward equation (differentiability of the stochastic flow in the
initial point; see \cite[Ch.~5]{Friedman} or \cite{Kunita}). This is quoted
and not proved here; the chapter reference is from memory. Conversely any
$u$ as in the theorem equals $\E\varphi(V_\tau^v)$, by It\^o's formula
applied to $u(\tau-\sigma,V_\sigma)$. Granting the quoted fact,
Theorem~\ref{thm:kin} gives: for every $\tau\ge0$, $v_{\tau/\eps^2}$
converges in law to $V_\tau^{v_0}$ as $\eps\to0$ (test functions in
$C_c^\infty\subset C_b^4$ determine convergence in law to a probability
measure), and for $\kappa=1$ the limit dynamics is reversible with respect
to the Maxwellian $M_T$ by Proposition~\ref{prop:einstein}.
\end{remark}

\begin{proof}[Proof of Theorem~\ref{thm:kin}]
Since $u(\tau)=\varphi+\int_0^\tau\cL_\kappa u\,\dd\tau'$ and
$\cL_\kappa u$ has two continuous bounded $v$-derivatives, one may
differentiate under the integral:
$\partial_\tau\partial_v^ju=\partial_v^j\cL_\kappa u$ for $|j|\le2$,
continuous and bounded by $C\mathcal{N}$. Fix $t$ and let
$u^0(\sigma,v)=u(\eps^2(t-\sigma),v)$ for $0\le\sigma\le t$, so
$\partial_\sigma u^0=-\eps^2\cL_\kappa u^0$. Let
\[
  \Theta_{ij}=Q\bigl(G_i\Psi_j-D_{ij}\bigr),\qquad
  \Xi_j=Q\bigl((G\cdot\nabla_v)\Psi_j-(\Div D)_j\bigr),
\]
both in $\cP_2$ by Lemma~\ref{lem:id}(c),(d) and Lemma~\ref{lem:poisson},
and (summing over repeated indices)
\[
  u^\eps=u^0+\eps\,\Psi_j\,\partial_ju^0
   +\eps^2\bigl(\Theta_{ij}\,\partial_{ij}u^0+\Xi_j\,\partial_ju^0\bigr).
\]
Compute $(\partial_\sigma+\cL)u^\eps$ with $\cL=A^v+\eps\cK$:
\begin{align*}
  \cL u^0&=\eps\,G_j\partial_ju^0,\\
  \cL(\eps\Psi_j\partial_ju^0)&=-\eps\,G_j\partial_ju^0
    +\eps^2\bigl[G_i\Psi_j\,\partial_{ij}u^0
      +(G\cdot\nabla_v)\Psi_j\,\partial_ju^0
      -\kappa T^{-1}(Dv)_j\partial_ju^0\bigr],\\
  A^v\bigl(\eps^2(\Theta_{ij}\partial_{ij}u^0+\Xi_j\partial_ju^0)\bigr)
   &=-\eps^2\bigl[(G_i\Psi_j-D_{ij})\partial_{ij}u^0
     +((G\cdot\nabla_v)\Psi_j-(\Div D)_j)\partial_ju^0\bigr].
\end{align*}
The terms of order $\eps$ cancel and those of order $\eps^2$ sum to
$\eps^2[D_{ij}\partial_{ij}+(\Div D)_j\partial_j-\kappa T^{-1}(Dv)_j\partial_j]u^0
=\eps^2\cL_\kappa u^0=-\partial_\sigma u^0$. What is left is
$(\partial_\sigma+\cL)u^\eps=\eps^3R_3+\eps^4R_4$,
\begin{align*}
  R_3&=-\Psi_j\,\partial_j\cL_\kappa u^0
     +\cK\bigl(\Theta_{ij}\,\partial_{ij}u^0+\Xi_j\,\partial_ju^0\bigr),\\
  R_4&=-\Theta_{ij}\,\partial_{ij}\cL_\kappa u^0-\Xi_j\,\partial_j\cL_\kappa u^0 ,
\end{align*}
where $\cK$ acts on the product (on the coefficient polynomials and, through
$G\cdot\nabla_v$, on the derivatives of $u^0$). $R_3$ involves derivatives
of $u$ up to order three and $R_4$ up to order four, so by
\eqref{eq:Kbound}
\[
  |R_3|\le C\mathcal{N}(1+|v|)(1+|z|)^3,\qquad
  |R_4|\le C\mathcal{N}(1+|z|)^2 .
\]
$u^\eps$ satisfies the hypotheses of Lemma~\ref{lem:dynkin}. At $\sigma=t$,
$u^0=\varphi$ and $|\E u^\eps(t)-\E\varphi(v_t)|\le C\mathcal{N}(\eps+\eps^2)$.
At $\sigma=0$, $|\E u^\eps(0)-u(\eps^2t,v_0)|\le C\mathcal{N}\eps^2$, because
$\E\Psi(v_0,z_0)=O(\eps)$ and $\Theta,\Xi$ are bounded in mean. By
Lemma~\ref{lem:diff},
$(\E(1+|v_\sigma|)^2)^{1/2}\le C(1+|v_0|)(1+\eps(1+\sigma)^{1/2})
\le C(1+|v_0|)(1+\tau^{1/2})$ for $\sigma\le t$. Hence
\[
  \bigl|\E\varphi(v_t)-u(\eps^2t,v_0)\bigr|
  \le C\mathcal{N}\eps+C\mathcal{N}(1+|v_0|)(1+\tau^{1/2})\,\eps^3t
     +C\mathcal{N}\eps^4t ,
\]
and $\eps^3t=\eps\tau$.
\end{proof}

\begin{remark}
Theorem~\ref{thm:kin} for $\kappa=0$ is a quantitative form of the
Kesten--Papanicolaou limit for this Markovian field; the content here is
the case $\kappa=1$, in which the limit acquires the friction
$-T^{-1}D(v)v$ and becomes the linear Fokker--Planck equation
$\partial_\tau f=\nabla\cdot(D(\nabla f+vf/T))$ with Maxwellian equilibrium.
The proof does not use ellipticity of $D$ (which may degenerate if $K$ does
not span $\mathbb{R}^d$), only the regularity of $u$.
\end{remark}

\section{Numerical check}

\begin{sloppypar}
The script \texttt{scripts/orbit\_backreaction\_check.py} (numpy only,
seed $20261124$) simulates \eqref{eq:kappa} for $d=2$, five modes,
$T=0.5$, $v_0=(0.8,0.3)$, $q=(1,0.5)$, dressed start; the parameters are in
the script and in the output file
\texttt{orbit\_backreaction\_check.json} in
\texttt{results/orbit\_backreaction/}.
\end{sloppypar}
\noindent For these parameters the formulas \eqref{eq:D}--\eqref{eq:m} give
\[
  \Div D(v_0)=(1.1884,\,0.3949),\qquad
  D(v_0)v_0/T=(1.3201,\,0.4124),\qquad q\cdot D(v_0)q=0.9941 .
\]
The modes are split exactly as $\tilde a=\hat a+\eps r$ with $\hat a$ a
stationary Ornstein--Uhlenbeck path, sampled exactly, and $r$ the solution
of a pathwise ordinary differential equation; $(x,v,r)$ are advanced by a
kick--drift--kick step of size $h=0.05$. The runs with $\kappa=0$ and
$\kappa=1$ share $\hat a$ (common random numbers), so the difference
isolates the back-reaction. For $\kappa=0$ the first Duhamel iterate, which
has mean exactly zero, is used as a control variate. The run took
$184$\,s on a machine that was heavily loaded by other jobs (an earlier
version with about a fifth less work took $123$\,s under lighter load). Three errors are kept apart: the
Monte Carlo standard error (s.e.), the time-step error (measured below),
and the theoretical error of the theorems.

\paragraph{(a) Mean velocity drift (Theorem~\ref{thm:drift}).} The table
gives $\E w_t/(\eps^2t)$, first component (the second behaves in the same
way and is in the output file). ``Paired'' is $\kappa=1$ minus $\kappa=0$
on the same noise, to be compared with $-D(v_0)v_0/T$.
\begin{center}\small\setlength{\tabcolsep}{4pt}
\begin{tabular}{rr|rr|rr}
  $\eps$ & $t$ & $\kappa=0$ (s.e.) & $\Div D$ & paired (s.e.) & $-Dv_0/T$ \\ \hline
  0.02 & 10 & 0.949 (0.018) & 1.188 & $-1.2890$ (0.0011) & $-1.3201$ \\
  0.02 & 20 & 1.047 (0.029) & 1.188 & $-1.3194$ (0.0030) & $-1.3201$ \\
  0.02 & 30 & 1.095 (0.031) & 1.188 & $-1.3425$ (0.0060) & $-1.3201$ \\
  0.02 & 40 & 1.083 (0.029) & 1.188 & $-1.3511$ (0.0097) & $-1.3201$ \\
  0.01 & 10 & 0.944 (0.027) & 1.188 & $-1.2840$ (0.0008) & $-1.3201$ \\
  0.01 & 20 & 1.046 (0.051) & 1.188 & $-1.3075$ (0.0021) & $-1.3201$ \\
  0.01 & 30 & 1.035 (0.064) & 1.188 & $-1.3105$ (0.0039) & $-1.3201$ \\
  0.01 & 40 & 1.069 (0.068) & 1.188 & $-1.3148$ (0.0062) & $-1.3201$ \\
\end{tabular}
\end{center}
($10^5$ samples at $\eps=0.02$, $5\times10^4$ at $\eps=0.01$.) The
deviations at $t=10$ are far outside the Monte Carlo error; they are the
theoretical error $C\eps^2$ of \eqref{eq:drift}, which is of relative size
$O(1/t)$. The increment between $t=10$ and $t=40$ removes that constant:
\begin{center}\small\setlength{\tabcolsep}{4pt}
\begin{tabular}{r|rr|rr}
  $\eps$ & $\kappa=0$ slope (s.e.) & $\Div D$ & paired slope (s.e.) & $-Dv_0/T$ \\ \hline
  0.02 & $(1.122,\,0.374)$ $(.036,\,.028)$ & $(1.188,\,0.395)$
       & $(-1.372,\,-0.435)$ $(.013,\,.008)$ & $(-1.320,\,-0.412)$ \\
  0.01 & $(1.130,\,0.390)$ $(.088,\,.066)$ & $(1.188,\,0.395)$
       & $(-1.325,\,-0.414)$ $(.008,\,.005)$ & $(-1.320,\,-0.412)$ \\
\end{tabular}
\end{center}
The polarization drag is confirmed to about $0.6\%$ at $\eps=0.01$ (within
one standard error). At $\eps=0.02$ the paired slope is $4\%$ off, four
standard errors; this is of the size of the relative theoretical error
$O(\eps^2t)$ expected from the variation of $m_\kappa$ along the diffusing
velocity (a heuristic second-order estimate
$\eps^2t\,m+\tfrac12\eps^4t^2\cL_\kappa m_\kappa$, not proved here, gives
$-1.398$ and $-1.339$ for the two rows), and it decreases with $\eps$ as it
should. The drift $\Div D$ without back-reaction is resolved less well,
because it is not helped by the pairing: the slopes agree with the formula
within $1.9$ standard errors at $\eps=0.02$ (a $3\%$ standard error), and at
$\eps=0.01$ the standard error is $8\%$, so that row is a weak test. The
sign and size of $\Div D$, and the fact that the generator is
$\nabla\cdot(D\nabla)$ rather than $D:\nabla^2$ (which would give zero
drift), are nevertheless clearly resolved at $\eps=0.02$.

\paragraph{(b) The $N^{1/3}$ law (Theorem~\ref{thm:third}).}
$4\times10^4$ samples; the s.e.\ of each characteristic function is $0.0034$
for $s=1$ and $0.0047$ for $s=1.5$.
\begin{center}\small\setlength{\tabcolsep}{4pt}
\begin{tabular}{rr|r|rr|rr}
  $N$ & $s$ & predicted & $\kappa=0$ & $\kappa=1$ & paired diff.\ (s.e.) & heur. \\ \hline
  $10^3$ & 1.0 & 0.7180 & $0.7244+0.0311i$ & $0.7308-0.0210i$ & $0.0064-0.0520i$ (0.0003) & $-0.0548i$ \\
  $10^3$ & 1.5 & 0.3268 & $0.3341+0.0092i$ & $0.3466-0.0452i$ & $0.0125-0.0544i$ (0.0009) & $-0.0561i$ \\
  $10^4$ & 1.0 & 0.7178 & $0.7250+0.0131i$ & $0.7265-0.0121i$ & $0.0016-0.0251i$ (0.0001) & $-0.0254i$ \\
  $10^4$ & 1.5 & 0.3274 & $0.3387+0.0004i$ & $0.3412-0.0260i$ & $0.0025-0.0265i$ (0.0004) & $-0.0261i$ \\
  $10^5$ & 1.0 & 0.7182 & $0.7211+0.0069i$ & $0.7214-0.0049i$ & $0.0004-0.0117i$ (0.0001) & $-0.0118i$ \\
  $10^5$ & 1.5 & 0.3272 & $0.3285-0.0000i$ & $0.3290-0.0121i$ & $0.0006-0.0121i$ (0.0002) & $-0.0121i$ \\
\end{tabular}
\end{center}
(The predicted value $\exp(-\eps^2t^3q\cdot Dq/3)$ varies slightly because
$t$ is rounded to the grid.) With and without back-reaction the real part
agrees with the same prediction: within one standard error at $N=10^5$,
and within three at $N=10^4$, where the deviations (about $0.01$) are
compatible with the proved $O(N^{-1/3})$ error ($N^{-1/3}=0.046$). The back-reaction is visible as an
imaginary part that decreases like $N^{-1/3}$. The ``heuristic'' column is
$-i\,e^{-s^3d_0/3}\,\eps^2t^2\,q\cdot D(v_0)v_0/2T$, the phase shift from
the friction displacement of Remark~\ref{rem:where}; it is not a proved
asymptotic for the paired difference, but agrees within about $5\%$. The
proved statement \eqref{eq:driftY} can be tested on the same runs through
$\E(Y^{\kappa=1}_t-Y^{\kappa=0}_t)$ against $-\eps^2t^2D(v_0)v_0/2T$: at
$N=10^5$, $s=1$, the measured value is
$(-0.014048,\,-0.004475)\pm(0.000012,\,0.000007)$ and the formula gives
$(-0.014211,\,-0.004440)$. The $1\%$ discrepancy is thirteen standard
errors and is theoretical error, of relative order $1/t=2\%$ by
\eqref{eq:driftY}.

\paragraph{(c) Time-step error.} On $2\times10^4$ samples at $\eps=0.02$,
the particle was advanced with steps $h$ and $2h$ on the same
Ornstein--Uhlenbeck paths. The differences in $\E w_t/(\eps^2t)$ between the
two step sizes were, at $t=20$ and $t=40$: for $\kappa=0$,
$-0.0001\pm0.0034$ and $-0.0007\pm0.0067$ (first component); for the paired
difference, $0.0006\pm0.0004$ and $0.0065\pm0.0070$. The time-step error is
therefore below the Monte Carlo error of the tables in (a), and not resolved.

\paragraph{Summary.} The numerical evidence supports the formula for the
polarization drag at the $1\%$ level, the divergence-form drift at the
$3$--$10\%$ level, and the unchanged $N^{1/3}$ law. It does not test the
rates in the error terms, nor Theorem~\ref{thm:kin}.

\section{What remains between this model and the particle system}

\begin{enumerate}
  \item \textbf{The medium.} The field modes relax at rates $\gamma_k$
    imposed by a thermostat. In the particle system the medium is the gas of
    other particles, the relaxation of its response is Landau damping, and
    the friction is the polarization drag
    computed with the dielectric function. The mechanism proved here (the
    drag is the response of the medium to the source $\propto k\cdot v$
    carried by the particle, and is tied to $D$ by the Gibbs structure) is
    the same, but the Markov property of the medium is not available there
    (gap~3 of the companion note).
  \item \textbf{Temperature.} The Einstein relation here is inherited from
    the thermostat at temperature $T$. For a non-Maxwellian background the
    Lenard--Balescu friction is $D$ applied to a different vector, and
    there is no Gibbs measure to guide the computation.
  \item \textbf{Sharp resonances.} All constants depend on
    $\min_k\gamma_k>0$ and on the finiteness of $K$.
  \item \textbf{Order $N^{-1/3}$.} The characteristic function of $Y$ is
    computed up to $O(N^{-1/3})$; the coefficient of that term is not
    (Remark~\ref{rem:next}).
  \item \textbf{Regularity input.} Theorem~\ref{thm:kin} is conditional on a
    classical regularity statement for the limiting Fokker--Planck equation
    (Remark~\ref{rem:reg}), and the invariance of the Gibbs measure on a
    quoted theorem (Remark~\ref{rem:inv}).
  \item \textbf{Self-consistency.} The tagged particle now acts on the
    field, but the field is not made of particles. Gap~1 of the companion
    note is untouched.
\end{enumerate}

\begin{thebibliography}{99}
\small
\bibitem{Balescu} R.~Balescu, Irreversible processes in ionized gases,
  \emph{Phys.\ Fluids} 3 (1960) 52--63.
\bibitem{CL} A.~O.~Caldeira and A.~J.~Leggett, Path integral approach to
  quantum Brownian motion, \emph{Physica A} 121 (1983) 587--616.
\bibitem{EK} S.~N.~Ethier and T.~G.~Kurtz, \emph{Markov Processes:
  Characterization and Convergence}, Wiley, 1986.
\bibitem{FKM} G.~W.~Ford, M.~Kac and P.~Mazur, Statistical mechanics of
  assemblies of coupled oscillators, \emph{J.\ Math.\ Phys.} 6 (1965)
  504--515.
\bibitem{Friedman} A.~Friedman, \emph{Stochastic Differential Equations and
  Applications}, vol.~1, Academic Press, 1975.
\bibitem{Hubbard} J.~Hubbard, The friction and diffusion coefficients of the
  Fokker--Planck equation in a plasma, \emph{Proc.\ R.\ Soc.\ A} 260 (1961)
  114--126.
\bibitem{KP} H.~Kesten and G.~C.~Papanicolaou, A limit theorem for
  stochastic acceleration, \emph{Comm.\ Math.\ Phys.} 78 (1980) 19--63.
\bibitem{Kubo} R.~Kubo, The fluctuation-dissipation theorem, \emph{Rep.\
  Prog.\ Phys.} 29 (1966) 255--284.
\bibitem{Kunita} H.~Kunita, \emph{Stochastic Flows and Stochastic
  Differential Equations}, Cambridge University Press, 1990.
\bibitem{Lenard} A.~Lenard, On Bogoliubov's kinetic equation for a spatially
  homogeneous plasma, \emph{Ann.\ Phys.} 10 (1960) 390--400.
\bibitem{PSV} G.~C.~Papanicolaou, D.~Stroock and S.~R.~S.~Varadhan,
  Martingale approach to some limit theorems, in \emph{Duke Turbulence
  Conference}, Duke Univ.\ Math.\ Ser.\ III, 1977.
\bibitem{PavliotisStuart} G.~A.~Pavliotis and A.~M.~Stuart, \emph{Multiscale
  Methods: Averaging and Homogenization}, Springer, 2008.
\bibitem{RR} N.~Rostoker and M.~N.~Rosenbluth, Test particles in a
  completely ionized plasma, \emph{Phys.\ Fluids} 3 (1960) 1--14.
\bibitem{SV} D.~W.~Stroock and S.~R.~S.~Varadhan, \emph{Multidimensional
  Diffusion Processes}, Springer, 1979.
\bibitem{TH} W.~B.~Thompson and J.~Hubbard, Long-range forces and the
  diffusion coefficients of a plasma, \emph{Rev.\ Mod.\ Phys.} 32 (1960)
  714--718.
\bibitem{Zwanzig} R.~Zwanzig, Nonlinear generalized Langevin equations,
  \emph{J.\ Stat.\ Phys.} 9 (1973) 215--220.
\end{thebibliography}
\noindent\emph{All bibliographic details above are from memory and were not
checked against the sources.}

\end{document}
